\section{A deterministic comparator for the same finite pricing target}

I compare the retained sampling procedures with exhaustive summation before
interpreting a sampling-query improvement as a financial pricing contribution.
This comparison uses the existing financial precision pool. It does not introduce
a newly optimised precision family or an additional statistical premise.
The result below is an elementary finite-sum identity, not a claim of theoretical
novelty. Its purpose is to identify a constructive competitor that the preceding
confidence budgets do not exclude.

\begin{proposition}[Exact finite summation and transfer to a continuous price]
Let the implemented financial output be
\(x:\mathcal D\longrightarrow\{0,1,\ldots,2^q\}\), where
\[
 \mathcal D=\prod_{j=1}^{d}\{0,1,\ldots,2^{b_j}-1\},
 \qquad N=2^{\sum_j b_j}.
\]
Suppose a deterministic certificate establishes
\[
 \left|P-\left(g_0+H\,\mathbb E[X]\right)\right|\leq\beta,
 \qquad X=2^{-q}x(U),
\]
for uniform \(U\) on \(\mathcal D\), a retained geometric correction centre
\(g_0\), and the continuous undiscounted price \(P\). The certificate includes
the error of \(g_0\). Then exhaustive evaluation gives the exact rational number
\[
 \overline x=\frac{\sum_{u\in\mathcal D}x(u)}{N2^q}
\]
and the deterministic price enclosure
\[
 P\in[g_0+H\overline x-\beta,\ g_0+H\overline x+\beta].
\]
There is no statistical error or random-bit independence assumption. If every
input state is evaluated exactly, the enclosure has confidence one under the
stated numerical and modelling premises.
\end{proposition}

\begin{proof}
The definition of a uniform random variable on a finite product gives
\(\mathbb E[X]=(N2^q)^{-1}\sum_{u\in\mathcal D}x(u)\). Every summand is an
integer. Arbitrary-precision integer summation therefore computes the numerator
without rounding. Substitution into the assumed deterministic certificate proves
both endpoints. No concentration inequality is used. In the retained one-input
enumeration, \(N\leq2^{17}\) and \(q\leq30\), so
\(\sum x(u)\leq N2^q\leq2^{47}<2^{64}\). An unsigned 64-bit sum is consequently
exact as well. This overflow argument concerns the output reduction; the
intermediate fixed-point products continue to use the separately proved limb
arithmetic of the financial interpreter.
\end{proof}

For direct, interval and full-conditioning graphs, respectively, the exhaustive
input sizes are
\[
 N_{\mathrm{direct}}=2^{d b_y+b_z},\qquad
 N_{\mathrm{interval}}=2^{d b_y+b_u},\qquad
 N_{\mathrm{full}}=2^{d b_y}.
\]
Thus storage for the scalar oracle tables, which is a sum of marginal table
sizes, does not determine the number of financial input combinations. These are
different quantities. A large tensor count is also not a lower bound on every
classical quadrature algorithm. I record every input size in the 273-prefix pool
and restrict actual exhaustive evaluation to the prespecified two-asset,
fully conditioned scope with at most \(2^{17}\) states.

If the prepared financial interpreter charges \(E\) scalar-lane operations per
input and uses \(P_v\) vector domain predicates, I retain the following
conservative online count for a one-dimensional enumeration:
\[
 C_{\mathrm{enum}}=NE+N+(N-1)+P_v(2N-1)+(N-1)+1
                =N(E+3+2P_v)-P_v-1.
\]
The terms charge financial arithmetic, input assignments, input increments,
domain comparisons and Boolean reductions, exact output addition, and final
rational construction. For the selected two-asset full graphs,
\(E=73\) and \(P_v=7\), giving \(90N-8\). These are algorithm-level scalar
operations, not processor instructions. The full table manifests, packed-word
reads, unpacked-word writes, graph-node counts and separate preparation times
are retained. Special-function setup is not converted into an invented processor
instruction count, and a shared cache is not interpreted as free preparation.

\section{A certified continuous two-asset comparator}

The following calculation also avoids the finite normal grid entirely. I retain
the frozen rational parameters and the same undiscounted price target. Write
\[
 S_i=K\exp(\mu_i+r_iY+\lambda Z),\qquad
 Y,Z\stackrel{\mathrm{ind}}{\sim}N(0,1),\qquad
 A=\sum_{i=1}^{2}w_iS_i,\qquad G=\prod_{i=1}^{2}S_i^{w_i},
\]
where \(w_i>0\), \(w_1+w_2=1\), \(K=100\), and \(\lambda>0\).
The retained exact rational check verifies
\(\mu_i+(r_i^2+\lambda^2)/2=(r-\delta)T\), so this is the stipulated
risk-neutral model. Define
\[
 m_A(y)=\log\sum_iw_i e^{\mu_i+r_iy},\qquad
 m_G(y)=\sum_iw_i(\mu_i+r_iy),\qquad
 F(s)=\min\{(s-K)^+,H\}.
\]
The common-factor conditional scalar price is
\(C_\lambda(m)=\mathbb E[F(K e^{m+\lambda Z})]\).

\begin{lemma}[Positive strike representation and analytic scalar kernel]
For real \(m\), let \(h=H/K\), \(a=\log(1+h)\),
\(t_0=m/\lambda\) and \(t_1=(m-a)/\lambda\). Then
\begin{align*}
 C_\lambda(m)
 &=K\int_1^{1+h}\Phi\!\left(\frac{m-\log k}{\lambda}\right)\,dk\\
 &=K\left\{e^{m+\lambda^2/2}
     [\Phi(t_0+\lambda)-\Phi(t_1+\lambda)]
       -[\Phi(t_0)-\Phi(t_1)]+h\Phi(t_1)\right\}.
\end{align*}
For positive real \(\lambda\), the displayed closed expression is an entire
function of the complex variable \(m\).
\end{lemma}

\begin{proof}
For every \(s\geq0\),
\(F(s)=\int_K^{K+H}\mathbf1_{\{s>x\}}\,dx\); the value at an endpoint is
irrelevant to this integral. Tonelli's theorem applies because the integrand is
nonnegative. Substitution \(x=Kk\) gives the first equality and the normal
probability \(\Phi((m-\log k)/\lambda)\).
For the second equality, integration by parts gives
\[
 \int_1^{1+h}\Phi\!\left(\frac{m-\log k}{\lambda}\right)\,dk
 =(1+h)\Phi(t_1)-\Phi(t_0)
 +\int_1^{1+h}\frac1\lambda
       \phi\!\left(\frac{m-\log k}{\lambda}\right)\,dk.
\]
In the last integral set \(k=e^{m-\lambda t}\). Completing the square yields
\[
 \int_{t_1}^{t_0}e^{m-\lambda t}\phi(t)\,dt
 =e^{m+\lambda^2/2}
       [\Phi(t_0+\lambda)-\Phi(t_1+\lambda)].
\]
Combining the boundary terms proves the formula. The normal CDF extension
\(\Phi(z)=(1+\operatorname{erf}(z/\sqrt2))/2\) is entire, as are the
exponential and their finite products and sums. The logarithm \(a\) is a
positive-real-parameter constant, not a logarithm of the variable \(m\).
\end{proof}

\begin{proposition}[Residual quadrature and a one-sided Gaussian tail enclosure]
Put \(\mu_G=\sum_iw_i\mu_i\), \(r_G=\sum_iw_ir_i\), and
\(\sigma_G=\sqrt{\lambda^2+r_G^2}\). For every \(L>0\),
\[
 P=\mathbb E[F(A)]
   =C_{\sigma_G}(\mu_G)
     +\int_{-L}^{L}
       [C_\lambda(m_A(y))-C_\lambda(m_G(y))]\phi(y)\,dy+\tau_L,
 \qquad 0\leq\tau_L\leq2H\Phi(-L).
\]
Suppose validated integration returns a real interval \([I_-,I_+]\) for the
finite integral, validated scalar evaluation returns \([g_-,g_+]\) for the
geometric correction, and \(T_+\geq2H\Phi(-L)\). Then
\[
 P\in[I_-+g_-,\ I_++g_++T_+].
\]
For the direct conditional arithmetic price, the analogous interval is obtained
by omitting the geometric correction and integrating
\(C_\lambda(m_A(y))\phi(y)\). Its omitted tail obeys the same bound.
\end{proposition}

\begin{proof}
Weighted arithmetic--geometric inequality gives \(A\geq G\) pointwise. Since
\(F\) is increasing and takes values in \([0,H]\),
\(0\leq F(A)-F(G)\leq H\). Conditioning on \(Y=y\) preserves both
inequalities and gives the residual kernel in the displayed integral.
Also \(\log(G/K)=\mu_G+r_GY+\lambda Z\), whose Gaussian variance is
\(r_G^2+\lambda^2\), proving the exact geometric correction.
The omitted normal probability is
\(\mathbb P(|Y|>L)=2\Phi(-L)\). Multiplication by the pointwise bound \(H\)
proves the one-sided tail enclosure. Addition of outward endpoints proves the
price interval. The direct case uses \(0\leq F(A)\leq H\) in precisely the
same tail argument.
\end{proof}

I use the existing authorised \texttt{python-flint} 0.9.0 implementation of
Arb/Acb validated ball arithmetic. The ball-enclosure foundation is established by \citet{johansson2017}; I consulted the author-hosted typeset version, Sections 2.2 and 5.1, and the current installed API contracts separately. The primary integral uses 128-bit precision,
\(L=8\), and a requested absolute integration tolerance \(10^{-16}\).
The requested tolerance is not used as an achieved error bound. The actual
returned finite ball, the separately returned correction ball and the outward
Gaussian-tail bound determine every endpoint. Each endpoint is exported as an
exact dyadic rational from its mantissa and exponent; decimal displays are
rounded outwards. The direct representation is separately computed with
176-bit precision, \(L=10\), and requested tolerance \(10^{-24}\).

The Acb quadrature routine evaluates complex neighbourhoods as well as real
points. Although the scalar cap kernel is entire, the log basket
\(m_A(y)=\log\sum_iw_i e^{\mu_i+r_iy}\) need not be analytic on an arbitrary
complex neighbourhood because its argument may vanish or meet the principal
logarithm's branch cut. I therefore forward the quadrature routine's analytic
verification flag to the complex logarithm. A non-finite analytic evaluation is
returned to the integrator, which must subdivide or fail to return a finite
enclosure. I do not replace this check with the fact that the integration path
is real. No clipping, midpoint substitution or heuristic quadrature error is
used to construct the price certificate. The assurance is conditional on the
validated-library enclosure contract and the checked mathematical integrand;
it is not a claim that this project proves the software implementation itself.

\section{Retained deterministic findings and their contribution boundary}

All 47 retained two-asset full-conditioning prefixes are exhaustively summed,
covering 330,752 outputs. A separately executed arbitrary-precision
Python-integer interpreter reproduces every output and every exact sum.
All deterministic radii are below \(0.005\), while the originally requested
undiscounted price error is \(0.01\). Every finite-grid certificate contains
the corresponding continuous Arb enclosure. The continuous calculation retains
all nine two-asset prices, with actual interval widths between approximately
\(1.2\times10^{-15}\) and \(3.14\times10^{-14}\). Every residual representation
overlaps its separately computed direct representation, and 126 scalar states
also pass independent positive strike-integral checks.

For these low-dimensional cases, a smaller quantum sampling-query count than a
Monte Carlo procedure is insufficient evidence of quantum pricing advantage.
Exact summation already requires fewer financial-output evaluations than the
retained certified Monte Carlo budget at every one of the 47 matched prefixes.
Validated continuous integration provides a still stronger constructive
competitor without marginal normal grids or a statistical failure allocation.
This is an adverse result for a low-dimensional quantum-advantage claim, not an
original quadrature theorem. Conditional smoothing and efficient deterministic
basket pricing already occur in the published literature, including
\citet{bayer}. I retain this finding to delimit where a wider contribution must
be sought. It does not prove a classical complexity lower bound, exclude
high-dimensional applications, establish superiority under a different financial
model, or convert host execution time into a physical quantum runtime.
